\documentclass[../main.tex]{subfiles}

\begin{document}
	\appendix
	\chapter{Metriche di valutazione}
	\label{app:metrics}
	
	Nel presente lavoro, la capacità predittiva e di generalizzazione dei modelli è stata quantificata utilizzando le seguenti metriche statistiche standard. Siano $y_i$ i valori reali (\textit{ground truth}), $\hat{y}_i$ i valori predetti dal modello, $\bar{y}$ la media dei valori reali e $N$ il numero totale dei campioni.
	
	\section*{Coefficiente di correlazione di Pearson ($R$)}
	Misura la correlazione lineare tra i valori reali e le predizioni, variando tra -1 e 1. Un valore prossimo a 1 indica una forte correlazione lineare positiva.
	\begin{equation}
		R = \frac{\sum_{i=1}^{N} (y_i - \bar{y})(\hat{y}_i - \bar{\hat{y}})}{\sqrt{\sum_{i=1}^{N} (y_i - \bar{y})^2 \sum_{i=1}^{N} (\hat{y}_i - \bar{\hat{y}})^2}}
	\end{equation}
	
	\section*{Coefficiente di determinazione ($R^2$)}
	Definito anche $R^2$, rappresenta la proporzione di varianza della variabile dipendente che viene spiegata dal modello. Nel contesto del machine learning, può assumere valori negativi se il modello risulta peggiore di una semplice predizione pari alla media dei dati costantemente.
	\begin{equation}
		R^2 = 1 - \frac{\sum_{i=1}^{N} (y_i - \hat{y}_i)^2}{\sum_{i=1}^{N} (y_i - \bar{y})^2}
	\end{equation}
	
	\section*{Root Mean Squared Error (RMSE)}
	L'errore quadratico medio sotto radice penalizza maggiormente gli scarti più grandi (outliers) rispetto alla semplice media, mantenendo la stessa unità di misura del parametro stimato.
	\begin{equation}
		\text{RMSE} = \sqrt{\frac{1}{N} \sum_{i=1}^{N} (y_i - \hat{y}_i)^2}
	\end{equation}
	
	\section*{Mean Absolute Error (MAE)}
	L'errore assoluto medio fornisce una misura lineare della distanza media tra le predizioni e i valori reali, ed è meno sensibile agli outliers rispetto al RMSE.
	\begin{equation}
		\text{MAE} = \frac{1}{N} \sum_{i=1}^{N} |y_i - \hat{y}_i|
	\end{equation}
	
	\section*{Median Absolute Error (MedAE)}
	L'errore assoluto mediano è il valore centrale degli scarti in valore assoluto. È una metrica estremamente robusta, in quanto ignora completamente la distorsione causata da singoli errori abnormi o anomalie nel dataset.
	\begin{equation}
		\text{MedAE} = \text{mediana}(|y_1 - \hat{y}_1|, |y_2 - \hat{y}_2|, \dots, |y_N - \hat{y}_N|)
	\end{equation}
	
	
	
	%B
	\chapter{Iperparametri per singolo target}
	\label{app:grafici_iperparametri}
	
	In questa appendice sono riportati i grafici di valutazione delle prestazioni predittive utilizzati per l'ottimizzazione degli iperparametri dei modelli lineari. A completamento dell'analisi condotta sulla viscosità globale (discussa nella Sezione \ref{sec:iperparametri}), vengono qui presentate le metriche disaggregate per i tre singoli parametri termodinamici dell'equazione MYEGA: la temperatura di transizione vetrosa ($T_g$), l'indice di fragilità ($m$) e la viscosità asintotica ad alta temperatura ($\log_{10}\eta_\infty$).
	
	
	\section*{Ridge Regression}
	
	Per offrire una visione chiara e compatta dell'effetto della regolarizzazione $L_2$, i risultati della regressione Ridge sono stati raggruppati per singolo target termodinamico. Per ciascun parametro vengono mostrati congiuntamente:
	\begin{itemize}
		\item L'andamento del coefficiente di correlazione di Pearson ($R$) e del coefficiente di determinazione ($R^2$) in funzione di $\alpha$, calcolati sia sui dati di \textit{training} che di \textit{test}.
		
		\item Le \textit{tracce Ridge}, ovvero l'evoluzione dei coefficienti di regressione normalizzati. La diminuzione della deviazione standard (bande sfumate, calcolate su 20 \textit{split} casuali) indica la stabilizzazione numerica del modello rispetto alla multicollinearità.
	\end{itemize}
	
	\begin{figure}[htbp]
		\centering
		\begin{subfigure}[b]{\textwidth}
			\centering
			\includegraphics[trim=0cm 0cm 0cm 0.8cm, clip, width=0.85\textwidth]{performance/ridge/Tg_fitted.pdf}
			\caption{Metriche di performance ($R$ e $R^2$).}
		\end{subfigure}
		
		\vspace{0.6cm}
		\begin{subfigure}[b]{\textwidth}
			\centering
			\includegraphics[trim=0cm 0cm 0cm 0.8cm, clip, width=0.85\textwidth]{tracce_ridge/Tg_fitted.pdf}
			\caption{Evoluzione dei pesi (tracce Ridge).}
		\end{subfigure}
		
		\caption{Regressione Ridge con target la temperatura di transizione vetrosa ($T_g$) al variare della penalità $\alpha$.}
		\label{fig:app_ridge_tg}
	\end{figure}
	
	\clearpage
	
	\begin{figure}[htbp]
		\centering
		\begin{subfigure}[b]{\textwidth}
			\centering
			\includegraphics[trim=0cm 0cm 0cm 0.8cm, clip, width=0.85\textwidth]{performance/ridge/m_fitted.pdf}
			\caption{Metriche di performance ($R$ e $R^2$).}
		\end{subfigure}
		
		\vspace{0.6cm}
		\begin{subfigure}[b]{\textwidth}
			\centering
			\includegraphics[trim=0cm 0cm 0cm 0.8cm, clip, width=0.85\textwidth]{tracce_ridge/m_fitted.pdf}
			\caption{Evoluzione dei pesi (tracce Ridge).}
		\end{subfigure}
		
		\caption{Regressione Ridge con target l'indice di fragilità ($m$) al variare della penalità $\alpha$.}
		\label{fig:app_ridge_m}
	\end{figure}
	
	\clearpage
	
	\begin{figure}[htbp]
		\centering
		\begin{subfigure}[b]{\textwidth}
			\centering
			\includegraphics[trim=0cm 0cm 0cm 0.8cm, clip, width=0.85\textwidth]{performance/ridge/log_eta_inf_fitted.pdf}
			\caption{Metriche di performance ($R$ e $R^2$).}
		\end{subfigure}
		
		\vspace{0.6cm}
		\begin{subfigure}[b]{\textwidth}
			\centering
			\includegraphics[trim=0cm 0cm 0cm 0.8cm, clip, width=0.85\textwidth]{tracce_ridge/log_eta_inf_fitted.pdf}
			\caption{Evoluzione dei pesi (tracce Ridge).}
		\end{subfigure}
		
		\caption{Regressione Ridge con target la viscosità asintotica ($\log_{10}\eta_\infty$) al variare della penalità $\alpha$.}
		\label{fig:app_ridge_eta}
	\end{figure}
	
	\clearpage
	
	\section*{Regressione Elastic Net}
	Le \textit{heatmap} illustrano i risultati della \textit{cross-validation} a 5 \textit{fold} sui dati di \textit{training}, mostrando la variazione di $R^2$ e RMSE al variare congiunto dell'intensità di regolarizzazione ($\alpha$) e del rapporto di miscelazione Lasso/Ridge ($\beta$).
	
	\begin{figure}[htbp]
		\centering
		\begin{subfigure}[b]{\textwidth}
			\centering
			\includegraphics[trim=0cm 0cm 0cm 1.5cm, clip, width=0.9\textwidth]{performance/elastic_net/ElasticNet_CV_Heatmap_Tg_fitted.pdf}
			\caption{Target: Temperatura di transizione vetrosa ($T_g$).}
		\end{subfigure}
		
		\vspace{0.3cm}
		\begin{subfigure}[b]{\textwidth}
			\centering
			\includegraphics[trim=0cm 0cm 0cm 1.5cm, clip, width=0.9\textwidth]{performance/elastic_net/ElasticNet_CV_Heatmap_m_fitted.pdf}
			\caption{Target: Indice di fragilità ($m$).}
		\end{subfigure}
		
		\vspace{0.3cm}
		\begin{subfigure}[b]{\textwidth}
			\centering
			\includegraphics[trim=0cm 0cm 0cm 1.5cm, clip, width=0.9\textwidth]{performance/elastic_net/ElasticNet_CV_Heatmap_log_eta_inf_fitted.pdf}
			\caption{Target: Viscosità asintotica ($\log_{10}\eta_\infty$).}
		\end{subfigure}
		
		\caption{\textit{Heatmap} delle prestazioni in \textit{cross-validation} per il modello Elastic Net sui singoli target.}
		\label{fig:app_elasticnet_targets}
		\vspace{-3cm}
	\end{figure}
	
	\clearpage
	
	\section*{Principal Component Regression (PCR)}
	Le curve mostrano la progressione dell'accuratezza predittiva all'aumentare del numero di componenti principali ($P$) incluse nel modello. In questo caso, le componenti sono aggiunte in ordine decrescente di varianza strutturale spiegata.
	
	\begin{figure}[htbp]
		\centering
		\begin{subfigure}[b]{\textwidth}
			\centering
			\includegraphics[trim=0cm 0cm 0cm 0.9cm, clip, width=0.8\textwidth]{performance/pcr/Tg_fitted_PCR_Performance.pdf}
			\caption{Target: Temperatura di transizione vetrosa ($T_g$).}
		\end{subfigure}
		
		\vspace{0.2cm}
		\begin{subfigure}[b]{\textwidth}
			\centering
			\includegraphics[trim=0cm 0cm 0cm 0.9cm, clip, width=0.8\textwidth]{performance/pcr/m_fitted_PCR_Performance.pdf}
			\caption{Target: Indice di fragilità ($m$).}
		\end{subfigure}
		
		\vspace{0.2cm}
		\begin{subfigure}[b]{\textwidth}
			\centering
			\includegraphics[trim=0cm 0cm 0cm 0.9cm, clip, width=0.8\textwidth]{performance/pcr/log_eta_inf_fitted_PCR_Performance.pdf}
			\caption{Target: Viscosità asintotica ($\log_{10}\eta_\infty$).}
		\end{subfigure}
		
		\caption{Metriche di performance del modello PCR al variare del numero di componenti principali per i singoli target.}
		\label{fig:app_pcr_targets}
		\vspace{-4cm}
	\end{figure}
	
	\clearpage

	\section*{Supervised PCR (SPCR)}
	
	A differenza della PCR standard, nei grafici seguenti le componenti principali sono inserite progressivamente nel modello ordinandole in base al valore assoluto della loro correlazione specifica con la grandezza target considerata.

	\begin{figure}[htbp]
		\centering
		\begin{subfigure}[b]{\textwidth}
			\centering
			\includegraphics[trim=0cm 0cm 0cm 0.9cm, clip, width=0.8\textwidth]{performance/spcr/Tg_fitted_SPCR_Performance.pdf}
			\caption{Target: Temperatura di transizione vetrosa ($T_g$).}
		\end{subfigure}
		
		\vspace{0.2cm}
		\begin{subfigure}[b]{\textwidth}
			\centering
			\includegraphics[trim=0cm 0cm 0cm 0.9cm, clip, width=0.8\textwidth]{performance/spcr/m_fitted_SPCR_Performance.pdf}
			\caption{Target: Indice di fragilità ($m$).}
		\end{subfigure}
		
		\vspace{0.2cm}
		\begin{subfigure}[b]{\textwidth}
			\centering
			\includegraphics[trim=0cm 0cm 0cm 0.9cm, clip, width=0.8\textwidth]{performance/spcr/log_eta_inf_fitted_SPCR_Performance.pdf}
			\caption{Target: Viscosità asintotica ($\log_{10}\eta_\infty$).}
		\end{subfigure}
		
		\caption{Metriche di performance del modello SPCR al variare del numero di componenti principali per i singoli target.}
		\label{fig:app_spcr_targets}
		\vspace{-4cm}
	\end{figure}
	
	
	
	
	%C
	\chapter{Analisi estesa delle predizioni}
	\label{app:grafici_performance_residui}
	
	In questa appendice sono raccolti i grafici di dettaglio relativi alla valutazione predittiva dei modelli analizzati. La prima parte presenta un confronto esteso delle distribuzioni dei residui e degli errori in funzione della distanza composizionale per tutte le architetture lineari testate e per la rete neurale ViscNet. La seconda parte è dedicata alla visualizzazione delle curve viscosimetriche per l'intero \textit{test set} (85 liquidi), includendo le bande di incertezza generate dal modello Ridge Regression.
	
	\section*{Qualità dell'interpolazione}
	
	I grafici seguenti mettono a confronto gli istogrammi 2D della densità predittiva per tutti i modelli valutati. La retta nera rappresenta la bisettrice ideale (predizione perfetta). Nei riquadri interni è riportata la distribuzione statistica dei residui.
	
	\begin{figure}[htbp]
		\centering
		\begin{subfigure}[b]{0.48\textwidth}
			\centering
			\includegraphics[width=\textwidth]{correlazione_viscnet.png}
			\caption{ViscNet \cite{cassarViscNetNeuralNetwork2021}}
		\end{subfigure}
		\hfill
		\begin{subfigure}[b]{0.48\textwidth}
			\centering
			\includegraphics[width=\textwidth]{correlazione/ridge_2D_density_histogram.png}
			\caption{Ridge Regression}
		\end{subfigure}
		
		\vspace{0.2cm}
		\begin{subfigure}[b]{0.48\textwidth}
			\centering
			\includegraphics[width=\textwidth]{correlazione/elastic_2D_density_histogram.png}
			\caption{Elastic Net}
		\end{subfigure}
		\hfill
		\begin{subfigure}[b]{0.48\textwidth}
			\centering
			\includegraphics[width=\textwidth]{correlazione/vif_ols_2D_density_histogram.png}
			\caption{VIF + OLS}
		\end{subfigure}
		
		\vspace{0.2cm}
		\begin{subfigure}[b]{0.48\textwidth}
			\centering
			\includegraphics[width=\textwidth]{correlazione/pcr_2D_density_histogram.png}
			\caption{PCR}
		\end{subfigure}
		\hfill
		\begin{subfigure}[b]{0.48\textwidth}
			\centering
			\includegraphics[width=\textwidth]{correlazione/spcr_2D_density_histogram.png}
			\caption{Supervised PCR}
		\end{subfigure}
		
		\caption{Istogrammi 2D della viscosità predittiva rispetto a quella sperimentale ($\log_{10}\eta$) per tutti i modelli analizzati. Risulta evidente il progressivo scostamento dalla retta ideale man mano che si impone maggiore sparsità ai modelli lineari.}
		\label{fig:app_2d_histograms}
	\end{figure}
	
	\clearpage
	
	\section*{Errore di generalizzazione}
	
	Di seguito è riportato l'andamento del RMSE in funzione della distanza chimica Canberra \cite{cassarViscNetNeuralNetwork2021} tra ciascun liquido del \textit{test set} e i campioni presenti nel \textit{training set}. L'analisi evidenzia il limite di estrapolazione al di fuori del dominio noto.
	
	\begin{figure}[htbp]
		\centering
		\begin{subfigure}[b]{0.4\textwidth}
			\centering
			\includegraphics[width=\textwidth]{distanza/viscnet_distance_plot.png}
			\caption{ViscNet \cite{cassarViscNetNeuralNetwork2021}}
		\end{subfigure}
		\hfill
		\begin{subfigure}[b]{0.4\textwidth}
			\centering
			\includegraphics[trim=0cm 0cm 0cm 0.9cm, clip, width=0.93\textwidth]{distanza/ridge_distance_plot.pdf}
			\caption{Ridge Regression}
		\end{subfigure}
		
		\begin{subfigure}[b]{0.4\textwidth}
			\centering
			\includegraphics[trim=0cm 0cm 0cm 0.9cm, clip, width=0.93\textwidth]{distanza/elastic_net_distance_plot.pdf}
			\caption{Elastic Net}
		\end{subfigure}
		\hfill
		\begin{subfigure}[b]{0.4\textwidth}
			\centering
			\includegraphics[trim=0cm 0cm 0cm 0.9cm, clip, width=0.93\textwidth]{distanza/vif_ols_distance_plot.pdf}
			\caption{VIF + OLS}
		\end{subfigure}
		
		\begin{subfigure}[b]{0.4\textwidth}
			\centering
			\includegraphics[trim=0cm 0cm 0cm 0.9cm, clip, width=0.93\textwidth]{distanza/pcr_distance_plot.pdf}
			\caption{PCR}
		\end{subfigure}
		\hfill
		\begin{subfigure}[b]{0.4\textwidth}
			\centering
			\includegraphics[trim=0cm 0cm 0cm 0.9cm, clip, width=0.93\textwidth]{distanza/spcr_distance_plot.pdf}
			\caption{Supervised PCR}
		\end{subfigure}
		
		\caption{Andamento del RMSE rispetto alla distanza chimica di Canberra dai campioni di addestramento. Si osserva un generale degrado delle prestazioni (aumento dell'errore) sui vetri con composizioni atipiche per tutti i modelli esaminati.}
		\label{fig:app_distance_plots}
		\vspace{-5cm}
	\end{figure}
	
	\clearpage
	
	\section*{Curve viscosimetriche sul test set (Ridge Regression)}
	
	Nelle pagine seguenti sono riportati i grafici viscosità-temperatura per tutti gli 85 liquidi che compongono l'\textit{unseen test set}. I punti (blu) rappresentano le misurazioni sperimentali estratte dal database originale.
	
	Le bande tratteggiate delineano l'intervallo di confidenza al 95\%, calcolato attraverso la propagazione dell'incertezza del modello su 100 iterazioni \textit{bootstrap}. Si noti come tali bande siano indicatrici dell'incertezza della regressione Ridge dovuta al campionamento e non del rumore intrinseco dei dati sperimentali.
	
	\vspace{1cm}
	
	\begin{figure}[htbp]
		\centering
		% Inserito il grafico singolo per ottimizzare lo spazio sotto il testo
		\includegraphics[height=5.5cm]{test_viscosity_plots/ridge/all_liquids_8.pdf}
		\caption{Misurazioni sperimentali e intervallo di confidenza al 95\% del modello Ridge Regression (1/85).}
		\label{fig:all_liquids_8}
	\end{figure}
	
	\clearpage
	
	\begin{figure}[htbp]
		\centering
		\includegraphics[width=\textwidth, height=\textheight, keepaspectratio]{test_viscosity_plots/ridge/all_liquids_1.pdf}
		\caption{Misurazioni sperimentali e intervallo di confidenza al 95\% del modello Ridge Regression (parte 1/7).}
		\label{fig:all_liquids_1}
	\end{figure}
	
	\clearpage
	
	\begin{figure}[htbp]
		\centering
		\includegraphics[width=\textwidth, height=\textheight, keepaspectratio]{test_viscosity_plots/ridge/all_liquids_2.pdf}
		\caption{Misurazioni sperimentali e intervallo di confidenza del modello (parte 2/7).}
		\label{fig:all_liquids_2}
	\end{figure}
	
	\clearpage
	
	\begin{figure}[htbp]
		\centering
		\includegraphics[width=\textwidth, height=\textheight, keepaspectratio]{test_viscosity_plots/ridge/all_liquids_3.pdf}
		\caption{Misurazioni sperimentali e intervallo di confidenza al 95\% del modello Ridge Regression (parte 3/7).}
		\label{fig:all_liquids_3}
	\end{figure}
	
	\clearpage
	
	\begin{figure}[htbp]
		\centering
		\includegraphics[width=\textwidth, height=\textheight, keepaspectratio]{test_viscosity_plots/ridge/all_liquids_4.pdf}
		\caption{Misurazioni sperimentali e intervallo di confidenza al 95\% del modello Ridge Regression (parte 4/7).}
		\label{fig:all_liquids_4}
	\end{figure}
	
	\clearpage
	
	\begin{figure}[htbp]
		\centering
		\includegraphics[width=\textwidth, height=0.9\textheight, keepaspectratio]{test_viscosity_plots/ridge/all_liquids_5.pdf}
		\caption{Misurazioni sperimentali e intervallo di confidenza al 95\% del modello Ridge Regression (parte 5/7).}
		\label{fig:all_liquids_5}
	\end{figure}
	
	\clearpage
	
	\begin{figure}[htbp]
		\centering
		\includegraphics[width=\textwidth, height=\textheight, keepaspectratio]{test_viscosity_plots/ridge/all_liquids_6.pdf}
		\caption{Misurazioni sperimentali e intervallo di confidenza al 95\% del modello Ridge Regression (parte 6/7).}
		\label{fig:all_liquids_6}
	\end{figure}
	
	\clearpage
	
	\begin{figure}[htbp]
		\centering
		\includegraphics[width=\textwidth, height=\textheight, keepaspectratio]{test_viscosity_plots/ridge/all_liquids_7.pdf}
		\caption{Misurazioni sperimentali e intervallo di confidenza al 95\% del modello Ridge Regression (parte 7/7).}
		\label{fig:all_liquids_7}
	\end{figure}
	
	
	%D
	\chapter{Features e Componenti Principali}
	\label{app:features}
	
	Di seguito sono riportate le prime 10 \textit{features} (ordinate in base al valore assoluto del peso assegnato) per ciascun approccio di regressione e per ogni \textit{target} ($T_g$, $m$, $\log_{10}\eta_\infty$).
	Assieme alle tabelle vengono presentati anche i grafici completi dell'importanza di tutte le \textit{features} (o componenti principali) per ciascun \textit{target}.
	
	Per attenuare le fluttuazioni legate al campionamento, i valori dei pesi riportati nelle tabelle e nei grafici sono stati ottenuti calcolando la media dei coefficienti estratti da 100 iterazioni indipendenti del modello. In ciascun ciclo, il dataset di \textit{training} è stato rimescolato e partizionato casualmente, assegnando il 50\% dei liquidi all'addestramento (\textit{training set}) e il restante 50\% alla validazione (\textit{test set}).
	
	\section*{Regressione Ridge}
	\vspace{-4mm}
	
	\begin{table}[htbp]
		\centering
		\begin{minipage}[c]{0.48\textwidth}
			\centering
			\resizebox{\linewidth}{!}{%
				\begin{tabular}{cl r}
					\toprule
					\textbf{Rank} & \textbf{Feature} & \textbf{Peso} \\
					\midrule
					1 & W\_Mean\_GSbandgap & $-0.220$ \\
					2 & W\_Std\_NsValence & $-0.183$ \\
					3 & W\_Sum\_GSbandgap & $-0.176$ \\
					4 & W\_Min\_NsValence & $-0.166$ \\
					5 & W\_Sum\_vdw\_radius\_uff & $0.161$ \\
					6 & W\_Sum\_nvalence & $0.160$ \\
					7 & W\_Sum\_NsUnfilled & $-0.159$ \\
					8 & W\_Sum\_NsValence & $0.159$ \\
					9 & W\_Sum\_atomic\_radius\_rahm & $0.151$ \\
					10 & W\_Mean\_FusionEnthalpy & $-0.150$ \\
					\bottomrule
				\end{tabular}%
			}
			\caption{Top 10 features per $T_g$ (Ridge).}
		\end{minipage}\hfill
		\begin{minipage}[c]{0.48\textwidth}
			\centering
			\includegraphics[trim=0.7cm 0cm 25cm 0.7cm, clip, width=\linewidth]{importanza_features/ridge_features_importance.pdf}
			\captionof{figure}{Importanza delle features per $T_g$.}
			\label{fig:ridge_features_tg}
			\vspace{-1.5cm}
		\end{minipage}
	\end{table}
	
	\begin{table}[htbp]
		\centering
		\begin{minipage}[c]{0.48\textwidth}
			\centering
			\resizebox{\linewidth}{!}{%
				\begin{tabular}{cl r}
					\toprule
					\textbf{Rank} & \textbf{Feature} & \textbf{Peso} \\
					\midrule
					1 & W\_Std\_num\_oxistates & $-0.183$ \\
					2 & A\_Std\_boiling\_point & $0.181$ \\
					3 & A\_Std\_FirstIonizationEnergy & $0.162$ \\
					4 & W\_Sum\_lattice\_constant & $0.162$ \\
					5 & W\_Min\_FirstIonizationEnergy & $0.155$ \\
					6 & A\_Max\_num\_oxistates & $0.155$ \\
					7 & W\_Mean\_NdValence & $0.154$ \\
					8 & A\_Min\_heat\_of\_formation & $0.146$ \\
					9 & W\_Min\_zeff & $-0.140$ \\
					10 & W\_Min\_ElectronAffinity & $-0.139$ \\
					\bottomrule
				\end{tabular}%
			}
			\caption{Top 10 features per $m$ (Ridge).}
		\end{minipage}\hfill
		\begin{minipage}[c]{0.48\textwidth}
			\centering
			\includegraphics[trim=13.5cm 0cm 12.5cm 0.7cm, clip, width=\linewidth]{importanza_features/ridge_features_importance.pdf}
			\captionof{figure}{Importanza delle features per $m$.}
		\end{minipage}
	\end{table}
	
	\begin{table}[htbp]
		\centering
		\begin{minipage}[c]{0.48\textwidth}
			\centering
			\resizebox{\linewidth}{!}{%
				\begin{tabular}{cl r}
					\toprule
					\textbf{Rank} & \textbf{Feature} & \textbf{Peso} \\
					\midrule
					1 & A\_Std\_atomic\_radius\_rahm & $-0.332$ \\
					2 & A\_Std\_FirstIonizationEnergy & $0.327$ \\
					3 & A\_Min\_ElectronAffinity & $-0.287$ \\
					4 & W\_Max\_dipole\_polarizability & $-0.260$ \\
					5 & W\_Min\_lattice\_constant & $-0.248$ \\
					6 & A\_Min\_FirstIonizationEnergy & $0.244$ \\
					7 & W\_Std\_NpUnfilled & $0.240$ \\
					8 & W\_Std\_SpaceGroupNumber & $-0.236$ \\
					9 & A\_Std\_NUnfilled & $0.236$ \\
					10 & A\_Std\_lattice\_constant & $-0.233$ \\
					\bottomrule
				\end{tabular}%
			}
			\caption{Top 10 features per $\log_{10}\eta_\infty$ (Ridge).}
		\end{minipage}\hfill
		\begin{minipage}[c]{0.48\textwidth}
			\centering
			\includegraphics[trim=26cm 0cm 0cm 0.7cm, clip, width=\linewidth]{importanza_features/ridge_features_importance.pdf}
			\captionof{figure}{Importanza delle features per $\log_{10}\eta_\infty$.}
		\end{minipage}
	\end{table}
	
	\clearpage
	\section*{Filtraggio VIF e Regressione OLS}
	
	\begin{table}[htbp]
		\centering
		\begin{minipage}[c]{0.48\textwidth}
			\centering
			\resizebox{\linewidth}{!}{%
				\begin{tabular}{cl r}
					\toprule
					\textbf{Rank} & \textbf{Feature} & \textbf{Peso} \\
					\midrule
					1 & W\_Max\_GSbandgap & $-0.270$ \\
					2 & W\_Sum\_vdw\_radius\_uff & $0.235$ \\
					3 & W\_Max\_atomic\_volume & $0.221$ \\
					4 & A\_Max\_NsUnfilled & $-0.191$ \\
					5 & A\_Max\_vdw\_radius\_uff & $0.170$ \\
					6 & A\_Max\_mass\_number & $-0.169$ \\
					7 & A\_Max\_FusionEnthalpy & $0.168$ \\
					8 & A\_Max\_NpUnfilled & $0.160$ \\
					9 & W\_Min\_NUnfilled & $-0.149$ \\
					10 & A\_Max\_zeff & $-0.147$ \\
					\bottomrule
				\end{tabular}%
			}
			\caption{Top 10 features per $T_g$ (VIF+OLS).}
		\end{minipage}\hfill
		\begin{minipage}[c]{0.48\textwidth}
			\centering
			\includegraphics[trim=0.7cm 0cm 25cm 0.7cm, clip, width=\linewidth]{importanza_features/vif_ols_features_importance.pdf}
			\captionof{figure}{Importanza delle features per $T_g$.}
		\end{minipage}
	\end{table}
	
	\begin{table}[htbp]
		\centering
		\begin{minipage}[c]{0.48\textwidth}
			\centering
			\resizebox{\linewidth}{!}{%
				\begin{tabular}{cl r}
					\toprule
					\textbf{Rank} & \textbf{Feature} & \textbf{Peso} \\
					\midrule
					1 & A\_Max\_FusionEnthalpy & $-0.259$ \\
					2 & W\_Max\_num\_oxistates & $-0.233$ \\
					3 & W\_Max\_c6\_gb & $0.163$ \\
					4 & A\_Max\_NsUnfilled & $-0.142$ \\
					5 & W\_Sum\_vdw\_radius\_uff & $-0.137$ \\
					6 & A\_Max\_mendeleev\_number & $0.133$ \\
					7 & A\_Max\_zeff & $0.132$ \\
					8 & A\_Mean\_num\_oxistates & $-0.117$ \\
					9 & W\_Std\_atomic\_volume & $-0.110$ \\
					10 & W\_Max\_SpaceGroupNumber & $-0.101$ \\
					\bottomrule
				\end{tabular}%
			}
			\caption{Top 10 features per $m$ (VIF+OLS).}
		\end{minipage}\hfill
		\begin{minipage}[c]{0.48\textwidth}
			\centering
			\includegraphics[trim=13.5cm 0cm 12.5cm 0.7cm, clip, width=\linewidth]{importanza_features/vif_ols_features_importance.pdf}
			\captionof{figure}{Importanza delle features per $m$.}
		\end{minipage}
		\vspace{-3cm}
	\end{table}
	
	\begin{table}[htbp]
		\centering
		\begin{minipage}[c]{0.48\textwidth}
			\centering
			\resizebox{\linewidth}{!}{%
				\begin{tabular}{cl r}
					\toprule
					\textbf{Rank} & \textbf{Feature} & \textbf{Peso} \\
					\midrule
					1 & W\_Sum\_vdw\_radius\_uff & $-0.292$ \\
					2 & A\_Max\_NsUnfilled & $0.190$ \\
					3 & A\_Max\_GSenergy\_pa & $-0.130$ \\
					4 & A\_Min\_SpaceGroupNumber & $0.129$ \\
					5 & W\_Min\_ElectronAffinity & $-0.114$ \\
					6 & A\_Min\_vdw\_radius\_uff & $0.106$ \\
					7 & A\_Sum\_GSmagmom & $-0.105$ \\
					8 & A\_Min\_max\_ionenergy & $-0.104$ \\
					9 & A\_Max\_mass\_number & $0.103$ \\
					10 & W\_Min\_c6\_gb & $-0.089$ \\
					\bottomrule
				\end{tabular}%
			}
			\caption{Top 10 features per $\log_{10}\eta_\infty$ (VIF+OLS).}
		\end{minipage}\hfill
		\begin{minipage}[c]{0.48\textwidth}
			\centering
			\includegraphics[trim=26cm 0cm 0cm 0.7cm, clip, width=\linewidth]{importanza_features/vif_ols_features_importance.pdf}
			\captionof{figure}{Importanza delle features per $\log_{10}\eta_\infty$.}
		\end{minipage}
		\vspace{1cm}
	\end{table}
	
	\newpage
	\section*{Regressione Elastic Net}
	
	\begin{table}[htbp]
		\centering
		\begin{minipage}[c]{0.48\textwidth}
			\centering
			\resizebox{\linewidth}{!}{%
				\begin{tabular}{cl r}
					\toprule
					\textbf{Rank} & \textbf{Feature} & \textbf{Peso} \\
					\midrule
					1 & W\_Sum\_GSbandgap & $-0.224$ \\
					2 & W\_Sum\_NsUnfilled & $-0.194$ \\
					3 & W\_Sum\_NsValence & $0.194$ \\
					4 & W\_Std\_density & $-0.102$ \\
					5 & W\_Max\_GSbandgap & $-0.088$ \\
					6 & A\_Std\_zeff & $-0.085$ \\
					7 & A\_Max\_vdw\_radius\_uff & $0.081$ \\
					8 & A\_Max\_zeff & $-0.071$ \\
					9 & A\_Max\_FusionEnthalpy & $0.064$ \\
					10 & A\_Max\_NpUnfilled & $0.063$ \\
					\bottomrule
				\end{tabular}%
			}
			\caption{Top 10 features per $T_g$ (Elastic Net).}
		\end{minipage}\hfill
		\begin{minipage}[c]{0.48\textwidth}
			\centering
			\includegraphics[trim=0.7cm 0cm 25cm 0.7cm, clip, width=\linewidth]{importanza_features/elastic_net_features_importance_no_zeroes.pdf}
			\captionof{figure}{Importanza delle features per $T_g$.}
		\end{minipage}
		\vspace{-3cm}
	\end{table}
	
	\begin{table}[htbp]
		\centering
		\begin{minipage}[c]{0.48\textwidth}
			\centering
			\resizebox{\linewidth}{!}{%
				\begin{tabular}{cl r}
					\toprule
					\textbf{Rank} & \textbf{Feature} & \textbf{Peso} \\
					\midrule
					1 & W\_Sum\_ElectronAffinity & $-0.232$ \\
					2 & A\_Max\_FusionEnthalpy & $-0.142$ \\
					3 & W\_Max\_SpaceGroupNumber & $-0.103$ \\
					4 & W\_Std\_num\_oxistates & $-0.096$ \\
					5 & W\_Sum\_num\_oxistates & $-0.090$ \\
					6 & W\_Sum\_NdValence & $0.079$ \\
					7 & A\_Min\_heat\_of\_formation & $0.071$ \\
					8 & W\_Std\_c6\_gb & $0.069$ \\
					9 & A\_Max\_ElectronAffinity & $0.063$ \\
					10 & A\_Max\_mendeleev\_number & $0.063$ \\
					\bottomrule
				\end{tabular}%
			}
			\caption{Top 10 features per $m$ (Elastic Net).}
		\end{minipage}\hfill
		\begin{minipage}[c]{0.48\textwidth}
			\centering
			\includegraphics[trim=13.5cm 0cm 12.5cm 0.7cm, clip, width=\linewidth]{importanza_features/elastic_net_features_importance_no_zeroes.pdf}
			\captionof{figure}{Importanza delle features per $m$.}
		\end{minipage}
	\end{table}
	
	\begin{table}[htbp]
		\centering
		\begin{minipage}[c]{0.48\textwidth}
			\centering
			\resizebox{\linewidth}{!}{%
				\begin{tabular}{cl r}
					\toprule
					\textbf{Rank} & \textbf{Feature} & \textbf{Peso} \\
					\midrule
					1 & W\_Max\_NsUnfilled & $0.125$ \\
					2 & W\_Sum\_vdw\_radius\_uff & $-0.089$ \\
					3 & A\_Max\_calc\_en\_sanderson & $-0.023$ \\
					4 & A\_Max\_vdw\_radius\_uff & $-0.019$ \\
					5 & W\_Min\_dipole\_polarizability & $-0.017$ \\
					6 & A\_Max\_GSmagmom & $-0.016$ \\
					7 & A\_Sum\_GSmagmom & $-0.016$ \\
					8 & A\_Std\_electrons & $0.014$ \\
					9 & A\_Std\_protons & $0.013$ \\
					10 & A\_Std\_atomic\_number & $0.013$ \\
					\bottomrule
				\end{tabular}%
			}
			\caption{Top 10 features per $\log_{10}\eta_\infty$ (Elastic Net).}
		\end{minipage}\hfill
		\begin{minipage}[c]{0.48\textwidth}
			\centering
			\includegraphics[trim=26cm 0cm 0cm 0.7cm, clip, width=\linewidth]{importanza_features/elastic_net_features_importance_no_zeroes.pdf}
			\captionof{figure}{Importanza delle features per $\log_{10}\eta_\infty$.}
		\end{minipage}
	\end{table}
	
	\clearpage
	\section*{Principal Component Regression (PCR)}
	
	\begin{table}[htbp]
		\centering
		\begin{tabular}{cl r r l r}
			\toprule
			\textbf{Rank} & \textbf{Comp.} & \textbf{Peso} & \textbf{Var. Spieg. (\%)} & \textbf{Top Feature} & \textbf{Load.} \\
			\midrule
			1 & PC6 & $-0.118$ & $4.01$ & A\_Std\_zeff & $0.126$ \\
			2 & PC14 & $-0.113$ & $1.00$ & W\_Std\_GSvolume\_pa & $0.150$ \\
			3 & PC13 & $0.078$ & $1.04$ & A\_Max\_nvalence & $0.178$ \\
			4 & PC8 & $0.071$ & $2.55$ & A\_Min\_mass\_number & $0.139$ \\
			5 & PC11 & $0.065$ & $1.59$ & A\_Max\_FusionEnthalpy & $0.194$ \\
			6 & PC5 & $-0.049$ & $6.89$ & A\_Std\_atomic\_radius & $0.144$ \\
			7 & PC7 & $-0.032$ & $2.77$ & A\_Mean\_NdUnfilled & $0.135$ \\
			8 & PC1 & $-0.027$ & $28.13$ & W\_Mean\_zeff & $0.080$ \\
			9 & PC12 & $-0.024$ & $1.29$ & A\_Mean\_GSmagmom & $0.220$ \\
			10 & PC4 & $0.015$ & $8.47$ & W\_Std\_NsValence & $0.107$ \\
			\bottomrule
		\end{tabular}
		\caption{Top 10 componenti principali per la predizione di $T_g$ (PCR).}
	\end{table}
	
	\begin{table}[htbp]
		\centering
		\begin{tabular}{cl r r l r}
			\toprule
			\textbf{Rank} & \textbf{Comp.} & \textbf{Peso} & \textbf{Var. Spieg. (\%)} & \textbf{Top Feature} & \textbf{Load.} \\
			\midrule
			1 & PC11 & $-0.162$ & $1.59$ & A\_Max\_FusionEnthalpy & $0.194$ \\
			2 & PC7 & $0.087$ & $2.77$ & A\_Mean\_NdUnfilled & $0.135$ \\
			3 & PC10 & $0.069$ & $1.86$ & W\_Std\_GSvolume\_pa & $0.127$ \\
			4 & PC12 & $-0.055$ & $1.29$ & A\_Mean\_GSmagmom & $0.220$ \\
			5 & PC6 & $-0.046$ & $4.01$ & A\_Std\_zeff & $0.126$ \\
			6 & PC4 & $-0.039$ & $8.47$ & W\_Std\_NsValence & $0.107$ \\
			7 & PC8 & $-0.034$ & $2.55$ & A\_Min\_mass\_number & $0.139$ \\
			8 & PC1 & $0.028$ & $28.13$ & W\_Mean\_zeff & $0.080$ \\
			9 & PC2 & $0.020$ & $16.98$ & A\_Mean\_atomic\_number & $0.089$ \\
			10 & PC14 & $-0.019$ & $1.00$ & W\_Std\_GSvolume\_pa & $0.150$ \\
			\bottomrule
		\end{tabular}
		\caption{Top 10 componenti principali per la predizione di $m$ (PCR).}
	\end{table}
	
	\begin{table}[htbp]
		\centering
		\begin{tabular}{cl r r l r}
			\toprule
			\textbf{Rank} & \textbf{Comp.} & \textbf{Peso} & \textbf{Var. Spieg. (\%)} & \textbf{Top Feature} & \textbf{Load.} \\
			\midrule
			1 & PC9 & $-0.053$ & $2.14$ & A\_Sum\_GSmagmom & $0.188$ \\
			2 & PC14 & $0.043$ & $1.00$ & W\_Std\_GSvolume\_pa & $0.150$ \\
			3 & PC12 & $0.029$ & $1.29$ & A\_Mean\_GSmagmom & $0.220$ \\
			4 & PC7 & $0.026$ & $2.77$ & A\_Mean\_NdUnfilled & $0.135$ \\
			5 & PC6 & $0.024$ & $4.01$ & A\_Std\_zeff & $0.126$ \\
			6 & PC3 & $0.018$ & $12.79$ & W\_Sum\_atomic\_radius & $0.117$ \\
			7 & PC4 & $-0.014$ & $8.47$ & W\_Std\_NsValence & $0.107$ \\
			8 & PC13 & $-0.013$ & $1.04$ & A\_Max\_nvalence & $0.178$ \\
			9 & PC11 & $-0.009$ & $1.59$ & A\_Max\_FusionEnthalpy & $0.194$ \\
			10 & PC2 & $-0.009$ & $16.98$ & A\_Mean\_atomic\_number & $0.089$ \\
			\bottomrule
		\end{tabular}
		\caption{Top 10 componenti principali per la predizione di $\log_{10}\eta_\infty$ (PCR).}
	\end{table}
	
	\begin{figure}[htbp]
		\centering
		\includegraphics[trim=0cm 0cm 0cm 0.7cm, clip, width=\textwidth]{importanza_features/pcr_features_importance.pdf}
		\vspace{0.5cm}
		\includegraphics[trim=0cm 0cm 0cm 0.7cm, clip, width=\textwidth]{importanza_features/pcr_projected_features_importance.pdf}
		\caption{In alto: Importanza delle Componenti Principali ($P=14$) per la predizione di $T_g$, $m$ e $\log_{10}\eta_\infty$ (PCR). In basso: Pesi delle proiezioni delle PC nello spazio delle originali \textit{features} chimico-fisiche.}
		\label{fig:pcr_plots}
	\end{figure}
	
	\clearpage
	\section*{Supervised PCR (SPCR)}
	
	\begin{table}[htbp]
		\centering
		\begin{tabular}{cl r r l r}
			\toprule
			\textbf{Rank} & \textbf{Comp.} & \textbf{Peso} & \textbf{Var. Spieg. (\%)} & \textbf{Top Feature} & \textbf{Load.} \\
			\midrule
			1 & PC46 & $0.224$ & $0.05$ & A\_Std\_vdw\_radius\_uff & $0.232$ \\
			2 & PC6 & $-0.118$ & $4.01$ & A\_Std\_zeff & $0.126$ \\
			3 & PC14 & $-0.113$ & $1.00$ & W\_Std\_GSvolume\_pa & $0.150$ \\
			4 & PC15 & $-0.091$ & $0.79$ & W\_Std\_NValence & $0.159$ \\
			5 & PC19 & $-0.088$ & $0.48$ & A\_Max\_mendeleev\_number & $0.416$ \\
			6 & PC13 & $0.078$ & $1.04$ & A\_Max\_nvalence & $0.178$ \\
			7 & PC8 & $0.071$ & $2.55$ & A\_Min\_mass\_number & $0.139$ \\
			8 & PC17 & $0.068$ & $0.71$ & A\_Min\_SpaceGroupNumber & $0.166$ \\
			9 & PC11 & $0.065$ & $1.59$ & A\_Max\_FusionEnthalpy & $0.194$ \\
			10 & PC5 & $-0.049$ & $6.89$ & A\_Std\_atomic\_radius & $0.144$ \\
			\bottomrule
		\end{tabular}
		\caption{Top 10 componenti principali per la predizione di $T_g$ (SPCR).}
	\end{table}
	
	\begin{table}[htbp]
		\centering
		\begin{tabular}{cl r r l r}
			\toprule
			\textbf{Rank} & \textbf{Comp.} & \textbf{Peso} & \textbf{Var. Spieg. (\%)} & \textbf{Top Feature} & \textbf{Load.} \\
			\midrule
			1 & PC53 & $-0.245$ & $0.03$ & W\_Min\_c6\_gb & $0.216$ \\
			2 & PC11 & $-0.162$ & $1.59$ & A\_Max\_FusionEnthalpy & $0.194$ \\
			3 & PC28 & $0.122$ & $0.18$ & A\_Max\_SpaceGroupNumber & $0.333$ \\
			4 & PC7 & $0.087$ & $2.77$ & A\_Mean\_NdUnfilled & $0.135$ \\
			5 & PC10 & $0.069$ & $1.86$ & W\_Std\_GSvolume\_pa & $0.127$ \\
			6 & PC12 & $-0.055$ & $1.29$ & A\_Mean\_GSmagmom & $0.220$ \\
			7 & PC6 & $-0.046$ & $4.01$ & A\_Std\_zeff & $0.126$ \\
			8 & PC4 & $-0.039$ & $8.47$ & W\_Std\_NsValence & $0.107$ \\
			9 & PC8 & $-0.034$ & $2.55$ & A\_Min\_mass\_number & $0.139$ \\
			10 & PC1 & $0.028$ & $28.13$ & W\_Mean\_zeff & $0.080$ \\
			\bottomrule
		\end{tabular}
		\caption{Top 10 componenti principali per la predizione di $m$ (SPCR).}
	\end{table}
	
	\begin{table}[htbp]
		\centering
		\begin{tabular}{cl r r l r}
			\toprule
			\textbf{Rank} & \textbf{Comp.} & \textbf{Peso} & \textbf{Var. Spieg. (\%)} & \textbf{Top Feature} & \textbf{Load.} \\
			\midrule
			1 & PC60 & $0.370$ & $0.02$ & A\_Mean\_GSmagmom & $0.164$ \\
			2 & PC59 & $-0.283$ & $0.02$ & W\_Max\_num\_oxistates & $0.240$ \\
			3 & PC46 & $-0.243$ & $0.05$ & A\_Std\_vdw\_radius\_uff & $0.232$ \\
			4 & PC34 & $0.197$ & $0.12$ & W\_Sum\_GSmagmom & $0.178$ \\
			5 & PC18 & $-0.072$ & $0.55$ & A\_Std\_calc\_en\_sanderson & $0.189$ \\
			6 & PC16 & $-0.069$ & $0.76$ & W\_Max\_atomic\_radius & $0.193$ \\
			7 & PC15 & $0.053$ & $0.79$ & W\_Std\_NValence & $0.159$ \\
			8 & PC9 & $-0.053$ & $2.14$ & A\_Sum\_GSmagmom & $0.188$ \\
			9 & PC14 & $0.043$ & $1.00$ & W\_Std\_GSvolume\_pa & $0.150$ \\
			10 & PC7 & $0.026$ & $2.77$ & A\_Mean\_NdUnfilled & $0.135$ \\
			\bottomrule
		\end{tabular}
		\caption{Top 10 componenti principali per la predizione di $\log_{10}\eta_\infty$ (SPCR).}
	\end{table}
	
	\begin{figure}[htbp]
		\centering
		\includegraphics[trim=0cm 0cm 0cm 0.7cm, clip, width=\textwidth]{importanza_features/spcr_features_importance.pdf}
		\vspace{0.5cm}
		\includegraphics[trim=0cm 0cm 0cm 0.7cm, clip, width=\textwidth]{importanza_features/spcr_projected_features_importance.pdf}
		\caption{In alto: Importanza delle Componenti Principali ($P=12$) per la predizione di $T_g$, $m$ e $\log_{10}\eta_\infty$ (Supervised PCR). In basso: Pesi delle proiezioni delle PC supervisionate nello spazio delle \textit{features} chimico-fisiche originali.}
		\label{fig:spcr_plots}
	\end{figure}
	
	
\end{document}